\documentclass[11pt, a4paper]{article}

\usepackage[T1]{fontenc}
\usepackage[utf8]{inputenc}
\usepackage[margin=1in]{geometry}
\usepackage{amsmath, amssymb}
\usepackage{graphicx}
\usepackage[hidelinks]{hyperref}
\usepackage{booktabs}

% The Constructions journal's own template loads a bibliography style
% called unified.bst, a linguistics-specific BibTeX style (the "Unified
% Style Sheet for Linguistics"). That exact file could not be obtained
% for this starter, and no BibTeX-compatible replacement ships in this
% TeX Live install, so natbib + plainnat is used instead — the standard,
% always-available substitute. See README.md.
\usepackage[round]{natbib}
\bibliographystyle{plainnat}

\title{A Constructions-Style Submission}
\author{Author Name \\ \small Institution}
\date{}

\begin{document}

\maketitle

\begin{abstract}
A one-paragraph abstract summarising the argument the paper makes,
written last but placed first.
\end{abstract}

\section{Introduction}

State the construction under analysis and the claim the paper makes
about it, with pointers to where each later section develops the
argument.

\section{Data and Method}

Describe the corpus or elicitation method used to gather the examples
discussed below. A numbered example in the standard linguistics format:

\begin{quote}
(1)\quad The cat sat on the mat.
\end{quote}

\section{Analysis}

Develop the argument, citing prior work as it becomes relevant
\citep{goldberg1995constructions}.

\begin{table}[ht]
\centering
\begin{tabular}{@{}ll@{}}
\toprule
Construction & Example \\
\midrule
Caused-motion & She sneezed the napkin off the table. \\
Ditransitive  & He gave her the book. \\
\bottomrule
\end{tabular}
\caption{Example constructions}
\label{tab:constructions}
\end{table}

\section{Conclusion}

Summarise the finding and what it implies for the broader constructionist
literature.

\bibliography{references}

\end{document}
